\section{Introduction}
\label{sec:introduction}

Organic semiconductors enable lightweight, large-area, and mechanically flexible
electronic devices. Their properties can be tuned through molecular design,
processing, film growth, and interfaces, which makes them attractive for displays,
sensors, low-cost circuits, and optoelectronic devices
\cite{Bruetting2005,Sirringhaus2005DevicePhysics}. This flexibility also makes
quantitative characterization difficult. Charge transport depends not only on the
molecular semiconductor, but also on grain boundaries, traps, dielectric
interfaces, contacts, and device geometry.

Small-molecule semiconductors such as C8-BTBT are especially useful model
systems for this question because they can form ordered thin films with clear
field-effect behaviour. At the same time, their measured device performance
remains sensitive to the complete stack, including the gate dielectric, the
metal contacts, and the processing history. A reliable characterization method
therefore has to distinguish between material transport in the accumulated film
and voltage losses or relaxation processes introduced by the device structure.

Gated thin-film test structures are widely used to evaluate such materials
because the gate voltage controls the charge density in a thin accumulation
layer. The commonly reported mobility, however, is not measured
directly. It is extracted from current--voltage data using a transport model and
therefore depends on the assumed device physics, the voltage range used for the
fit, and nonideal contributions. In a conventional two-terminal transistor measurement, the measured voltage
includes both the channel voltage drop and voltage losses at the metal--organic
contacts. Gate-dependent injection and contact resistance can
therefore distort transfer curves and lead to apparent mobilities that do not
represent the semiconductor sheet alone \cite{Natali2012ChargeInjection,
Bittle2016MobilityOverestimation}.

The gated van der Pauw method addresses this problem by combining field-effect
modulation with a four-terminal sheet-conductance measurement. A current is
driven through one pair of contacts, while a separate pair of voltage probes
measures the transverse voltage. Under suitable conditions, contact voltage drops
at the current-injecting electrodes are not included directly in the extracted
sheet resistance. The method can therefore provide a mobility that is less
sensitive to contact losses, although it still requires stable current injection,
low leakage, negligible voltage-probe current, and a well-defined accumulation
regime.

The aim of this thesis is to characterize gate-induced electrical transport in
C8-BTBT organic thin films using the current-driven gated van der Pauw method.
Conventional two-terminal field-effect measurements are used as complementary
reference measurements to verify gate modulation and to identify
contact-sensitive device behaviour. The central analysis determines the sheet conductance,
long-range thin-film mobility, and threshold voltage of the accumulated
semiconductor sheet from four-terminal measurements. Measurements of the
complete current-injection path and of the transient circuit response are used
to examine the remaining limitations of the characterization. The central
question is therefore whether the electrical properties of the accumulated
organic thin film can be determined without directly including voltage drops at
the current-injecting contacts.